\documentclass[11pt]{article}
\usepackage[margin=1in]{geometry}
\usepackage[hidelinks]{hyperref}
\usepackage{enumitem}
\usepackage{fancyvrb}
\usepackage{xcolor}

\setlength{\parindent}{0pt}
\setlist[itemize]{noitemsep, topsep=2pt}
\setlist[enumerate]{noitemsep, topsep=2pt}

\title{RDKDC ROS 2 Jazzy Docker Setup}
\author{EN.530.646 Robot Devices, Kinematics, Dynamic and Control}
\date{}

\begin{document}
\maketitle

\section*{Meeting Notes}

\begin{enumerate}
    \item We moved the course setup toward ROS 2 Jazzy in Docker so students do not need to install Ubuntu or ROS directly on their laptops.
    \item We found that host MATLAB does not reliably communicate with ROS 2 inside Docker Desktop through raw DDS networking on Windows. To make the setup robust, MATLAB talks to a local HTTP bridge, and the bridge talks to ROS 2 inside Docker.
    \item We verified the full prototype on Windows: Docker/VNC started, the bridge worked, a ROS 2 launch flow brought up UR5e/RViz/controllers, and MATLAB called \texttt{ur5\_interface()} to move the simulated UR5e to a test pose and return it home.
\end{enumerate}

\section*{What Students Install}

Students install only the following:

\begin{itemize}
    \item Docker Desktop
    \item MATLAB
    \item MATLAB ROS Toolbox
    \item MATLAB Robotics System Toolbox
    \item The course repository
\end{itemize}

Students do not install Ubuntu, ROS 2, Gazebo, RViz, or Universal Robots packages directly on the host machine. Those are provided inside Docker.

\section*{How The Setup Works}

The runtime layout is:

\begin{Verbatim}
MATLAB on Windows/macOS
    -> http://127.0.0.1:8765
    -> RDKDC HTTP bridge inside Docker
    -> ROS 2 Jazzy, RViz, Gazebo, UR5e simulation inside Docker
\end{Verbatim}

The bridge is necessary because Docker Desktop networking can block or confuse normal ROS 2 DDS discovery between host MATLAB and the Linux container.

\section*{Generated Files}

The start script creates or updates these files automatically.

On Windows:

\begin{Verbatim}
C:\Users\<student>\ros2_ws\rdkdc_bridge\rdkdc_http_bridge.py
C:\Users\<student>\ros2_ws\matlab\ur5_interface.m
C:\Users\<student>\ros2_ws\matlab\ur5e_interface.m
C:\Users\<student>\ros2_ws\matlab\tf_frame.m
C:\Users\<student>\ros2_ws\src\rdkdc_setup\launch\ur5e_sim.launch.py
\end{Verbatim}

On macOS:

\begin{Verbatim}
~/ros2_ws/rdkdc_bridge/rdkdc_http_bridge.py
~/ros2_ws/matlab/ur5_interface.m
~/ros2_ws/matlab/ur5e_interface.m
~/ros2_ws/matlab/tf_frame.m
~/ros2_ws/src/rdkdc_setup/launch/ur5e_sim.launch.py
\end{Verbatim}

If \texttt{ros2\_ws} already exists, the setup does not delete the workspace. It creates the needed folders if missing and overwrites only the generated bridge/helper/setup files listed above.

\section*{Windows Setup}

Open PowerShell in the course repository folder and run:

\begin{Verbatim}
Set-ExecutionPolicy -Scope Process -ExecutionPolicy Bypass
.\build_ros_windows.ps1
.\start_ros_windows.ps1
\end{Verbatim}

The build script builds the Docker image. The start script starts Docker, mounts the student's local workspace into Docker, starts noVNC, starts the MATLAB bridge, and builds the generated ROS 2 launch package.

The workspace mount is:

\begin{Verbatim}
Host:      C:\Users\<student>\ros2_ws
Docker:    /root/ros2_ws
\end{Verbatim}

Any files edited under the host \texttt{ros2\_ws} are immediately visible inside Docker at \texttt{/root/ros2\_ws}.

The start script will:

\begin{enumerate}
    \item Start the Docker container.
    \item Start noVNC at \url{http://127.0.0.1:6080}.
    \item Start the RDKDC HTTP bridge at \url{http://127.0.0.1:8765}.
    \item Generate and build the \texttt{rdkdc\_setup} ROS 2 launch package in \texttt{/root/ros2\_ws}.
    \item Run a MATLAB bridge smoke test if MATLAB is found.
\end{enumerate}

To stop:

\begin{Verbatim}
.\stop_ros_windows.ps1
\end{Verbatim}

\section*{macOS Setup}

Open Terminal in the course repository folder and run:

\begin{Verbatim}
chmod +x build_ros_mac.sh start_ros_mac.sh stop_ros_mac.sh
./build_ros_mac.sh
./start_ros_mac.sh
\end{Verbatim}

The workspace mount is:

\begin{Verbatim}
Host:      ~/ros2_ws
Docker:    /root/ros2_ws
\end{Verbatim}

Any files edited under the host \texttt{ros2\_ws} are immediately visible inside Docker at \texttt{/root/ros2\_ws}.

To stop:

\begin{Verbatim}
./stop_ros_mac.sh
\end{Verbatim}

\section*{Launching UR5e and RViz}

After the start script finishes, open the browser VNC desktop. In the Linux terminal inside the Docker/VNC desktop, run:

\begin{Verbatim}
source /root/ros2_ws/install/setup.bash
ros2 launch rdkdc_setup ur5e_sim.launch.py ur_type:=ur5e
\end{Verbatim}

This launch file starts the UR5e simulation, starts RViz, and activates the controllers needed by \texttt{ur5\_interface()}.

Useful options:

\begin{Verbatim}
ros2 launch rdkdc_setup ur5e_sim.launch.py ur_type:=ur5e
ros2 launch rdkdc_setup ur5e_sim.launch.py ur_type:=ur5
ros2 launch rdkdc_setup ur5e_sim.launch.py launch_rviz:=false
ros2 launch rdkdc_setup ur5e_sim.launch.py gazebo_gui:=true
\end{Verbatim}

\section*{Student MATLAB Workflow}

After setup finishes and the launch file is running inside Docker, students should write normal MATLAB lab code. They should not need to write bridge setup code in each assignment.

Example:

\begin{Verbatim}
ur5e = ur5_interface();

q = [0.40; -1.20; 0.80; -1.90; 0.30; 0.20];
ur5e.move_joints(q, 4);
pause(5);

ur5e.move_joints(ur5e.home, 4);
\end{Verbatim}

If MATLAB cannot find \texttt{ur5\_interface}, add the generated MATLAB folder once:

Windows:
\begin{Verbatim}
addpath('C:\Users\<student>\ros2_ws\matlab')
\end{Verbatim}

macOS:
\begin{Verbatim}
addpath('~/ros2_ws/matlab')
\end{Verbatim}

\section*{Verify The Setup}

The browser should open:

\begin{Verbatim}
http://127.0.0.1:6080/vnc.html?autoconnect=true&resize=scale
\end{Verbatim}

Inside the VNC desktop, RViz should show the UR5e after running \texttt{ros2 launch rdkdc\_setup ur5e\_sim.launch.py}. In MATLAB, a simple motion test should move the simulated robot:

\begin{Verbatim}
ur5e = ur5_interface();
disp(ur5e.get_current_joints().')
ur5e.move_joints([0.40; -1.20; 0.80; -1.90; 0.30; 0.20], 4);
pause(5);
disp(ur5e.get_current_joints().')
ur5e.move_joints(ur5e.home, 4);
\end{Verbatim}

\section*{Ports}

The setup uses these local ports:

\begin{itemize}
    \item \texttt{6080}: browser VNC
    \item \texttt{5901}: native VNC
    \item \texttt{8765}: MATLAB HTTP bridge
    \item \texttt{11811/udp}: Fast DDS discovery port, kept available for diagnostics
\end{itemize}

The scripts remove old course containers that commonly hold these ports. If an unrelated application is already using one of these ports, Docker will report a port binding error. Close the application or ask a TA to help identify the process.

\section*{Stopping The Environment}

Windows PowerShell:

\begin{Verbatim}
.\stop_ros_windows.ps1
\end{Verbatim}

macOS Terminal:

\begin{Verbatim}
./stop_ros_mac.sh
\end{Verbatim}

\section*{TA Notes}

\begin{itemize}
    \item Students should use the build/start/stop scripts for their OS. The start script prepares Docker, VNC, the bridge, and the launch package. Students then run \texttt{ros2 launch rdkdc\_setup ur5e\_sim.launch.py ur\_type:=ur5e} or \texttt{ur\_type:=ur5} inside Docker for simulation.
    \item The bridge-backed MATLAB files are generated into the student's home workspace, not the course repo root.
    \item Avoid shipping another \texttt{ur5\_interface.m} in the current MATLAB working directory unless it is the bridge-backed version, because MATLAB path precedence can load the wrong file.
\end{itemize}

\end{document}
